\section{Unified Interpretation}
    \label{sec:unification}
    
    \subsection{Overview}
        
        Swirl-String Theory (SST) proposes a unified interpretation in which fundamental physical phenomena emerge from the dynamics of a continuous medium supporting stable vortex structures.
        
        This section synthesizes the results of the previous sections into a coherent conceptual framework.
    
    \subsection{Ontology of the Theory}
        
        \textbf{[CORE CLAIM]}:
        The fundamental ontology consists of:
        
        \begin{itemize}
            \item A continuous incompressible medium (Section~\ref{sec:axioms})
            \item Stable vortex filament configurations (Section~\ref{sec:geometry})
            \item Circulation as a conserved topological invariant
        \end{itemize}
        
        \textbf{Interpretation:}
        Particles are not elementary objects, but persistent dynamical structures.
    
    \subsection{Mass as Trapped Circulation Energy}
        
        From Section~\ref{sec:foundations} and Section~\ref{sec:delay}:
        
        \begin{itemize}
            \item Energy is stored in rotational motion
            \item Circulation defines the energy scale
        \end{itemize}
        
        \begin{align}
            E = m c^2
        \end{align}
        
        \textbf{[UNIFIED INTERPRETATION]}:
        \begin{center}
            Mass corresponds to localized, self-sustained circulation energy in the medium.
        \end{center}
    
    \subsection{Time as a Relational Quantity}
        
        From Section~\ref{sec:time} and Section~\ref{sec:foundations}:
        
        \begin{itemize}
            \item Time is measured through physical processes
            \item The Swirl-Clock encodes local temporal rates
        \end{itemize}
        
        \begin{align}
            S_{(t)} = \sqrt{1 - \frac{v^2}{c^2}}
        \end{align}
        
        \textbf{[UNIFIED INTERPRETATION]}:
        \begin{center}
            Time is not fundamental, but emerges from the kinematic state of the medium.
        \end{center}
    
    \subsection{Discreteness without Quantization Postulates}
        
        From Section~\ref{sec:delay}:
        
        \begin{itemize}
            \item Delay introduces temporal nonlocality
            \item Stability selects discrete modes
        \end{itemize}
        
        \textbf{[UNIFIED INTERPRETATION]}:
        \begin{center}
            Discrete physical spectra arise dynamically from delay-induced mode selection.
        \end{center}
        
        This replaces operator-based quantization with a dynamical mechanism.
    
    \subsection{Charge and Circulation}
        
        \textbf{[SPECULATIVE]}:
        
        Charge-like properties may be associated with circulation orientation and magnitude:
        
        \begin{itemize}
            \item Sign of circulation $\leftrightarrow$ charge sign
            \item Circulation magnitude $\leftrightarrow$ charge strength
        \end{itemize}
    
    \subsection{Unification of Interactions}
        
        Within SST:
        
        \begin{itemize}
            \item Electromagnetic effects arise from transverse torsion modes
            \item Gravitational effects arise from large-scale circulation geometry
            \item Inertial mass arises from localized rotational energy
        \end{itemize}
        
        \textbf{[SPECULATIVE]}:
        These interactions are not independent forces but different manifestations of the same underlying fluid dynamics.
    
    \subsection{Atomic Structure as an Emergent Layer}
        
        From Section~\ref{sec:atomic}:
        
        \begin{itemize}
            \item Standard quantum mechanics appears as an effective description
            \item SST introduces circulation-dependent corrections
        \end{itemize}
        
        \textbf{[UNIFIED INTERPRETATION]}:
        \begin{center}
            Quantum mechanics is an emergent effective theory of underlying vortex dynamics.
        \end{center}
    
    \subsection{Consistency with Spectroscopy}
        
        From Section~\ref{sec:spectroscopy}:
        
        \begin{itemize}
            \item Internal modes are suppressed by an excitation gap
            \item Low-energy physics remains unchanged
        \end{itemize}
        
        \textbf{[CRITICAL]}:
        The unified picture does not violate known experimental constraints.
    
    \subsection{Hierarchy of Description}
        
        The theory admits multiple levels:
        
        \begin{enumerate}
            \item \textbf{Microscopic level:} vortex filament dynamics
            \item \textbf{Mesoscopic level:} delay-induced mode structure
            \item \textbf{Macroscopic level:} effective quantum description
        \end{enumerate}
    
    \subsection{Conceptual Summary}
        
        The SST framework can be summarized as:
        
        \begin{center}
            \begin{tabular}{ll}
                \textbf{Mass} & $\rightarrow$ trapped circulation energy \\
                \textbf{Time} & $\rightarrow$ relational swirl-clock field \\
                \textbf{Charge} & $\rightarrow$ circulation property \\
                \textbf{Quantization} & $\rightarrow$ delay-induced stability \\
            \end{tabular}
        \end{center}
    
    \subsection{Canonical Statement}
        
        \begin{center}
            \textit{
                All observed physical structure emerges from the dynamics of a continuous medium in which stable vortex configurations, governed by quantized circulation and delay dynamics, produce the phenomena conventionally attributed to particles, fields, and quantum mechanics.
            }
        \end{center}
    
    \subsection{Scope and Limits}
        
        \textbf{[IMPORTANT]}:
        
        \begin{itemize}
            \item The framework is not yet a complete fundamental theory
            \item Several mechanisms remain to be derived rigorously
            \item Many results are conditional or phenomenological
        \end{itemize}
    
    \subsection{Final Consistency Principle}
        
        \begin{center}
            \textit{
                A valid unified interpretation must preserve all experimentally verified results while providing a deeper structural explanation of their origin.
            }
        \end{center}